% generator_I29.tex -- Prop I.29 (Theor.): parallel lines => alternate angles equal; also external = internal and co-interior = 2 right angles.
\documentclass[12pt]{article}\usepackage{geometry}\geometry{a5paper,margin=1.5cm}
\usepackage[lining]{ebgaramond}\usepackage{amssymb}\usepackage{byrne}
\def\mpPre{u := 1cm; textLabels := false;}
\begin{document}
\defineNewPicture[1/5]{%
  pair A,B,C,D,E,F,G,H,I,J,d[];%
  A:=(0,0);B:=A shifted (7/2u,0);d1:=(0,-2u);%
  C:=A shifted d1;D:=B shifted d1;E:=11/20[A,B];F:=9/20[C,D];%
  G:=7/4[F,E];H:=4/3[E,F];d2:=(3/2u,1/2u);%
  I:=E shifted -d2;J:=E shifted d2;%
  byAngleDefine(I,E,A,byblue,SOLID_SECTOR);%
  byAngleDefine(H,E,I,byyellow,SOLID_SECTOR);%
  byAngleDefine(B,E,H,byblack,SOLID_SECTOR);%
  byAngleDefine(G,E,B,byred,SOLID_SECTOR);%
  byAngleDefine(G,F,D,byblack,SOLID_SECTOR);%
  draw byNamedAngleResized();%
  draw byLine(I,E,byblack,SOLID_LINE,REGULAR_WIDTH);%
  draw byLine(E,J,byblack,DASHED_LINE,REGULAR_WIDTH);%
  draw byLine(A,B,byyellow,SOLID_LINE,REGULAR_WIDTH);%
  draw byLine(C,D,byred,SOLID_LINE,REGULAR_WIDTH);%
  draw byLine(G,H,byblue,SOLID_LINE,REGULAR_WIDTH);%
  draw byLabelLine(0)(AB,CD,GH);%
  draw byLabelsOnPolygon(A,E,G)(OMIT_FIRST_LABEL+OMIT_LAST_LABEL,0);%
  draw byLabelsOnPolygon(H,F,C)(OMIT_FIRST_LABEL+OMIT_LAST_LABEL,0);%
  draw byLabelPoint(I,lineAngle.IE-90,1);%
  draw byLabelPoint(J,lineAngle.EJ+90,1);%
}
\noindent\begin{minipage}[t]{0.45\linewidth}\centering\vspace{0pt}%
\drawCurrentPicture\end{minipage}\hfill
\begin{minipage}[t]{0.52\linewidth}\vspace{0pt}%
\textbf{Book~I.\enspace Prop.\ XXIX. Theor.}
\medskip\noindent\itshape
A straight line (\drawUnitLine{GH}) falling on two parallel
straight lines (\drawUnitLine{AB} and \drawUnitLine{CD}), makes
the alternate angles equal to one another; and also the
external equal to the internal and opposite angle on the same
side; and the two internal angles on the same side together
equal to two right angles.
\bigskip\noindent\upshape
For if the alternate angles \drawAngle{IEA,HEI} and
\drawAngle{GFD} be not equal, draw \drawUnitLine{IE}, making
$\drawAngle{HEI}=\drawAngle{GFD}$~(pr.~I.23).\\[4pt]
Therefore $\drawUnitLine{IE,EJ}\parallel\drawUnitLine{CD}$~(pr.~I.27)
and therefore two straight lines which intersect are parallel
to the same straight line, which is impossible~(ax.~XII).\\[4pt]
Hence \drawAngle{IEA,HEI} and \drawAngle{GFD} are not unequal,
that is, they are equal:
$\drawAngle{IEA,HEI}=\drawAngle{GEB}$~(pr.~I.15);
$\therefore \drawAngle{GEB}=\drawAngle{GFD}$, the external
angle equal to the internal and opposite on the same side:
if \drawAngle{BEH} be added to both, then
$\drawAngle{GFD}+\drawAngle{BEH}=\drawAngle{BEH,GEB}=\drawTwoRightAngles$~(pr.~I.13).
That is to say, the two internal angles at the same side of
the cutting line are equal to two right angles.
\begin{flushright}Q.\ E.\ D.\end{flushright}
\end{minipage}
\end{document}
